%%=============================================================================
%% Inleiding
%%=============================================================================

\chapter{\IfLanguageName{dutch}{Inleiding}{Introduction}}%
\label{ch:inleiding}


De sector van financiële technologie, of Fintech, heeft de afgelopen vijf jaar een jaarlijkse groei van 23,58\% doorgemaakt \autocite{statista2024fintechgrowth}. Een belangrijk deel van die groei is toe te schrijven aan de inzet van autonome AI-agents: softwaresystemen die zelfstandig beslissingen nemen en acties uitvoeren \autocite{russell2021artificialintelligence}. Ze worden niet alleen voor eenvoudige taken gebruikt, maar ook voor complexe bedrijfsprocessen zoals het afhandelen van klantvragen of het verwerken van transacties. Om te communiceren met interne databases en API's is het Model Context Protocol (MCP) de voorbije jaren de standaard geworden. MCP biedt een uniforme interface waardoor AI-agents platformonafhankelijk met bedrijfssystemen kunnen communiceren, vergelijkbaar met de rol die HTTP speelde voor webapplicaties \autocite{anthropic2024mcpwhitepaper}. Zonder zo'n gestandaardiseerd protocol zou elke organisatie voor elk AI-model een aparte integratie moeten bouwen, met alle gevolgen voor onderhoud en schaalbaarheid van dien.

Die standaardisatie heeft echter ook een keerzijde op vlak van beveiliging. Neem een Fintech-bedrijf dat via MCP een AI-agent toegang geeft tot zijn betaalsystemen. Aanvallers kunnen in zo'n opstelling \emph{prompt injection} gebruiken om de instructies van de agent te manipuleren door gewijzigde invoer in het systeem te injecteren. Via \emph{tool poisoning} kunnen zij schadelijke commando's verbergen in schijnbaar legitieme datastromen. En via \emph{privilege escalation} kan een agent door een beveiligingsfout onbedoeld verhoogde toegangsrechten verkrijgen, waardoor kritieke systemen bereikbaar worden voor onbevoegden \autocite{song2025protocolunveilingattackvectors}.

Security-teams hebben doorgaans te weinig middelen om MCP-implementaties grondig te beveiligen tegen deze specifieke dreigingen. De huidige praktijk is afhankelijk van handmatige code-reviews en ad-hoc beveiligingstools die niet specifiek zijn ontworpen voor de technische eigenheid van het Model Context Protocol. Deze algemene beveiligingstools (zoals Burp Suite voor API-testing of Semgrep voor code-analyse) missen de nodige contextuele kennis om MCP-specifieke kwetsbaarheden, zoals gemanipuleerde tool-\hspace{0pt}beschrijvingen of onveilige JSON-RPC-communicatie, systematisch te detecteren. Dit gebrek aan gespecialiseerde tooling creëert een bottleneck in het ontwikkelproces: ofwel vertraagt innovatie door langdurige handmatige beveiligingscontroles, ofwel worden MCP-servers in productie gebracht zonder grondige veiligheidsbeoordeling. Dat 56\% van de ondernemingen die AI-agents inzetten beveiligingsincidenten rapporteert in hun MCP-omgevingen maakt de urgentie concreet \autocite{guo2025systematicanalysismcpsecurity}. Dit onderzoek bouwt een security assessment framework dat MCP-servers automatisch en systematisch test.


%De inleiding moet de lezer net genoeg informatie verschaffen om het onderwerp te begrijpen en in te zien waarom de onderzoeksvraag de moeite waard is om te onderzoeken. In de inleiding ga je literatuurverwijzingen beperken, zodat de tekst vlot leesbaar blijft. Je kan de inleiding verder onderverdelen in secties als dit de tekst verduidelijkt. Zaken die aan bod kunnen komen in de inleiding~\autocite{Pollefliet2011}:

%\begin{itemize}
%  \item context, achtergrond
%  \item afbakenen van het onderwerp
%  \item verantwoording van het onderwerp, methodologie
%  \item probleemstelling
%  \item onderzoeksdoelstelling
%  \item onderzoeksvraag
%  \item \ldots
%\end{itemize}

\section{\IfLanguageName{dutch}{Probleemstelling}{Problem Statement}}%
\label{sec:probleemstelling}

Fintech-bedrijven beschikken momenteel niet over een gestructureerde en geautomatiseerde manier om MCP-servers op kritieke veiligheidskwetsbaarheden te controleren vóór live-gang. Security teams vertrouwen op handmatige code reviews, een aanpak die arbeidsintensief en foutgevoelig is. Ad-hoc tools zoals Burp Suite of Semgrep bieden slechts gedeeltelijke dekking omdat ze de specifieke context van MCP-aanvallen missen.

\begin{itemize}
  \item \textbf{Specifieke risico's:} Zonder degelijke screening zijn servers kwetsbaar voor:
        \begin{itemize}
          \item \emph{Prompt injection attacks} via manipulatie van tool-beschrijvingen.
          \item \emph{Token theft \& privilege escalation}, waarbij agents te brede rechten krijgen.
          \item \emph{Tool poisoning}, waarbij kwaadaardige instructies in data worden verstopt.
          \item \emph{Confused deputy attacks}, omdat servers vaak geen onderscheid maken tussen gebruikers.
          \item Gebrek aan \emph{audit logging}, waardoor acties niet traceerbaar zijn.
        \end{itemize}
  \item \textbf{Doelgroep:} Het onderzoek richt zich op security teams, DevSecOps-engineers en enterprise architects binnen Fintech, e-commerce en andere gereguleerde industrieën.
\end{itemize}

%Uit je probleemstelling moet duidelijk zijn dat je onderzoek een meerwaarde heeft voor een concrete doelgroep. De doelgroep moet goed gedefinieerd en afgelijnd zijn. Doelgroepen als ``bedrijven,'' ``KMO's'', systeembeheerders, enz.~zijn nog te vaag. Als je een lijstje kan maken van de personen/organisaties die een meerwaarde zullen vinden in deze bachelorproef (dit is eigenlijk je steekproefkader), dan is dat een indicatie dat de doelgroep goed gedefinieerd is. Dit kan een enkel bedrijf zijn of zelfs één persoon (je co-promotor/opdrachtgever).

\section{\IfLanguageName{dutch}{Onderzoeksvraag}{Research question}}%
\label{sec:onderzoeksvraag}

\paragraph{Hoofdvraag}
Hoe kan een framework worden ontwikkeld dat MCP-server implementaties automatisch test op een vooraf gedefinieerde set kritieke kwetsbaarheden, en hoe verhoudt de dekkingsgraad en benodigde tijd per assessment zich tot een handmatige beoordeling op basis van twee concrete evaluatiecases?

\paragraph{Deelvragen Probleemdomein}
\begin{enumerate}
  \item Welke kwetsbaarheden uit de OWASP MCP Top 10 zijn het meest kritiek voor Fintech-toepassingen, en hoe worden ze geclassificeerd in de bestaande literatuur?
  \item Welke beperkingen hebben generieke beveiligingstools zoals Burp Suite en Semgrep bij het detecteren van MCP-specifieke kwetsbaarheden?
  \item Hoeveel tijd vereist een handmatige beoordeling van een MCP-server implementatie volgens beschikbare richtlijnen, en wat zijn de voornaamste bronnen van foutmarge?
\end{enumerate}

\paragraph{Deelvragen Oplossingsdomein}
\begin{enumerate}
  \item Welk percentage van de ingebedde kwetsbaarheden in de evaluatiecases detecteert het framework correct (recall), en wat is de bijhorende false positive rate?
  \item Hoeveel tijd heeft het framework nodig voor een volledige assessment in vergelijking met een gestructureerde handmatige procedure op dezelfde target?
  \item Welke categorieën kwetsbaarheden detecteert het framework betrouwbaarder dan handmatige review, en in welke gevallen blijft menselijke expertise noodzakelijk?
\end{enumerate}

%Wees zo concreet mogelijk bij het formuleren van je onderzoeksvraag. Een onderzoeksvraag is trouwens iets waar nog niemand op dit moment een antwoord heeft (voor zover je kan nagaan). Het opzoeken van bestaande informatie (bv. ``welke tools bestaan er voor deze toepassing?'') is dus geen onderzoeksvraag. Je kan de onderzoeksvraag verder specifiëren in deelvragen. Bv.~als je onderzoek gaat over performantiemetingen, dan

\section{\IfLanguageName{dutch}{Onderzoeksdoelstelling}{Research objective}}%
\label{sec:onderzoeksdoelstelling}

Deze bachelorproef beoogt een geautomatiseerd security assessment framework te ontwerpen, ontwikkelen en evalueren dat specifiek is toegespitst op MCP-server implementaties. De evaluatie steunt op twee concrete cases: een opzettelijk kwetsbare testserver (\emph{Goat}) en een representatieve open-source MCP-server. De concrete deliverables zijn:
\begin{enumerate}
  \item Een werkend prototype (Proof-of-Concept) van de security scanner die kwetsbaarheden uit de OWASP MCP Top 10 kan detecteren via statische analyse en dynamische fuzzing.
  \item Een kwantitatieve evaluatie van de effectiviteit van het framework op basis van twee cases, gemeten aan de hand van:
        \begin{itemize}
          \item Detectiegraad (recall): het percentage correct geïdentificeerde kwetsbaarheden
          \item Tijdsverschil: benodigde tijd voor een volledige assessment tegenover een gestructureerde handmatige procedure
          \item False positive rate: nauwkeurigheid van de detectie
        \end{itemize}
  \item Een set evidence-based aanbevelingen voor de veilige implementatie van MCP, onderbouwd door de bevindingen uit de evaluatiefase, met expliciete vermelding van de beperkingen van de externe validiteit.
\end{enumerate}


%Wat is het beoogde resultaat van je bachelorproef? Wat zijn de criteria voor succes? Beschrijf die zo concreet mogelijk. Gaat het bv.\ om een proof-of-concept, een prototype, een verslag met aanbevelingen, een vergelijkende studie, enz.

\section{\IfLanguageName{dutch}{Opzet van deze bachelorproef}{Structure of this bachelor thesis}}%
\label{sec:opzet-bachelorproef}

% Het is gebruikelijk aan het einde van de inleiding een overzicht te
% geven van de opbouw van de rest van de tekst. Deze sectie bevat al een aanzet
% die je kan aanvullen/aanpassen in functie van je eigen tekst.

De rest van deze bachelorproef is als volgt opgebouwd:

In Hoofdstuk~\ref{ch:stand-van-zaken} wordt een overzicht gegeven van de stand van zaken binnen het onderzoeksdomein, op basis van een literatuurstudie.

In Hoofdstuk~\ref{ch:methodologie} wordt de methodologie toegelicht en worden de gebruikte onderzoekstechnieken besproken om een antwoord te kunnen formuleren op de onderzoeksvragen.

% TODO: Vul hier aan voor je eigen hoofstukken, één of twee zinnen per hoofdstuk

In Hoofdstuk~\ref{ch:conclusie}, tenslotte, wordt de conclusie gegeven en een antwoord geformuleerd op de onderzoeksvragen. Daarbij wordt ook een aanzet gegeven voor toekomstig onderzoek binnen dit domein.